% cox and mullender, semaphores.
% 
% process has one thread with two stacks
%  
% pike et al, sleep and wakeup
\chapter{Escalonamento}
\label{CH:SCHED}

Qualquer sistema operacional provavelmente será executado com mais processos do que a quantidade de CPUs disponíveis no computador, então é necessário um plano para compartilhar o tempo das CPUs entre os processos. Idealmente, esse compartilhamento seria transparente para os processos de usuário. Uma abordagem comum é fornecer a cada processo a ilusão de que possui sua própria CPU virtual, por meio da
\indextext{multiplexação}
dos processos sobre as CPUs físicas.
Este capítulo explica como o xv6 realiza essa multiplexação.

%%
\section{Multiplexação}
%%

O xv6 realiza multiplexação trocando cada CPU de um processo para outro em duas situações. Primeiro, o mecanismo de
\lstinline{sleep}
e
\lstinline{wakeup}
do xv6 realiza a troca quando um processo faz uma chamada de sistema que bloqueia (ou seja, que precisa esperar por um evento), geralmente em
\lstinline{read},
\lstinline{wait},
ou
\lstinline{sleep}.
Segundo, o xv6 força periodicamente uma troca para lidar com processos que computam por longos períodos sem bloquear.
As primeiras são trocas voluntárias;
as segundas são chamadas de involuntárias.
Essa multiplexação cria a ilusão de que cada processo possui sua própria CPU.

Implementar essa multiplexação impõe alguns desafios. Primeiro, como alternar de um processo para outro?
A ideia básica é salvar e restaurar os registradores da CPU, embora o fato de que isso não pode ser expresso em C torne a tarefa complicada.
Segundo, como forçar trocas de forma transparente para os processos de usuário? O xv6 usa a técnica padrão na qual interrupções de um temporizador de hardware acionam trocas de contexto.
Terceiro, todas as CPUs alternam entre o mesmo conjunto de processos, então é necessário um plano de bloqueios (locking) para evitar condições de corrida.
Quarto, a memória de um processo e outros recursos devem ser liberados quando o processo termina, mas ele mesmo não pode fazer tudo isso — por exemplo, ele não pode liberar sua própria pilha do kernel enquanto ainda a está utilizando.
Quinto, cada CPU de uma máquina multi-core precisa lembrar qual processo está executando, para que as chamadas de sistema afetem o estado correto do processo no kernel.
Finalmente,
\lstinline{sleep}
e
\lstinline{wakeup}
permitem que um processo libere a CPU e espere ser acordado por outro processo ou por uma interrupção.
É necessário tomar cuidado para evitar condições de corrida que resultem na perda de notificações de wakeup.
%%
\section{Código: Troca de contexto}
%%

\begin{figure}[t] \center \includegraphics[scale=0.5]{fig/switch.pdf} \caption{Troca de um processo de usuário para outro. Neste exemplo, o xv6 está rodando com uma CPU (e, portanto, uma thread de escalonador).} \label{fig:switch} \end{figure}

A Figura~\ref{fig:switch}
resume os passos envolvidos na troca de um processo de usuário para outro:
um trap (chamada de sistema ou interrupção) do espaço do usuário para a thread do kernel do processo antigo,
uma troca de contexto para a thread do escalonador da CPU atual,
uma troca de contexto para a thread do kernel de um novo processo,
e um retorno do trap para o processo no espaço do usuário.
O xv6 tem threads separadas
(registradores e pilhas salvos)
nas quais o escalonador é executado porque
não é seguro que o escalonador execute sobre
a pilha do kernel de algum processo: outra CPU pode
acordar o processo e executá-lo, e seria um desastre
usar a mesma pilha em duas CPUs diferentes.
Existe uma thread de escalonador separada para cada CPU, para lidar
com situações nas quais mais de uma CPU está executando
um processo que deseja liberar a CPU.

%
% Acho que nunca é seguro. Se o processo antigo está como ZOMBIE, o
% processo pai pode marcá-lo como UNUSED e outro núcleo pode alocá-lo em
% fork(). Se o processo antigo está como SLEEPING ou RUNNABLE, outro núcleo pode
% acordá-lo e executá-lo.
%
% O que muda é que, se excluíssemos o wait(), seria
% seguro que um processo em exit() alterasse seu estado diretamente para UNUSED
% (sem ZOMBIE, sem necessidade de limpeza pelo pai).
%

Nesta seção, vamos examinar os mecanismos de troca
entre uma thread do kernel e uma thread do escalonador.

Trocar de uma thread para outra envolve salvar os registradores da CPU da thread antiga
e restaurar os registradores previamente salvos da
nova thread; o fato de que
o ponteiro da pilha e o contador de programa
são salvos e restaurados significa que a CPU trocará de pilha e
também mudará o código que está executando.

A função
\indexcode{swtch}
salva e restaura registradores durante uma troca de thread no kernel.
\lstinline{swtch}
não conhece diretamente as threads; ela apenas salva e
restaura conjuntos de registradores RISC-V, chamados
\indextext{contexts}.
Quando é hora de um processo liberar a CPU,
a thread do kernel do processo chama
\lstinline{swtch}
para salvar seu próprio contexto e restaurar o contexto do escalonador.
Cada contexto está contido em uma
\lstinline{struct}
\lstinline{context}
\lineref{kernel/proc.h:/^struct.context/},
que por sua vez está contida em uma
\lstinline{struct proc}
do processo ou em uma
\lstinline{struct cpu}
da CPU.
\lstinline{swtch}
recebe dois argumentos:
\indexcode{struct context}
\lstinline{*old}
e
\lstinline{struct}
\lstinline{context}
\lstinline{*new}.
Ela salva os registradores atuais em
\lstinline{old},
carrega os registradores de
\lstinline{new},
e retorna.

Vamos seguir um processo ao longo de
\lstinline{swtch}
até o escalonador.
Vimos no Capítulo~\ref{CH:TRAP}
que uma das possibilidades ao final de uma interrupção é que
\indexcode{usertrap}
chame
\indexcode{yield}.
\lstinline{yield}
por sua vez chama
\indexcode{sched},
que chama
\indexcode{swtch}
para salvar o contexto atual em
\lstinline{p->context}
e alternar para o contexto do escalonador previamente salvo em
\indexcode{cpu->context}
\lineref{kernel/proc.c:/swtch..p/}.

\lstinline{swtch}
\lineref{kernel/swtch.S:/swtch/}
salva apenas os registradores que devem ser preservados pelo callee (callee-saved);
o compilador C gera código no chamador para salvar
os registradores que devem ser preservados pelo chamador (caller-saved) na pilha.
\lstinline{swtch} conhece o deslocamento (offset) de cada
campo de registrador em
\lstinline{struct context}.
Ela não salva o contador de programa.
Em vez disso,
\lstinline{swtch}
salva o registrador
\lstinline{ra},
que contém o endereço de retorno de onde
\lstinline{swtch}
foi chamada.
Agora
\lstinline{swtch}
restaura os registradores a partir do novo contexto,
que contém os valores de registradores salvos por uma
\lstinline{swtch}
anterior.
Quando
\lstinline{swtch}
retorna, ela volta para as instruções apontadas pelo registrador
\lstinline{ra}
restaurado, ou seja,
a instrução de onde a nova thread havia chamado
\lstinline{swtch}
anteriormente.
Além disso, ela retorna usando a pilha da nova thread,
já que é para lá que o registrador restaurado
\lstinline{sp}
aponta.

No nosso exemplo,
\indexcode{sched}
chamou
\indexcode{swtch}
para alternar para
\indexcode{cpu->context},
o contexto do escalonador específico daquela CPU.
Esse contexto foi salvo em um momento anterior, quando
\lstinline{scheduler}
chamou
\lstinline{swtch}
\lineref{kernel/proc.c:/swtch.&c/}
para alternar para o processo que agora está liberando a CPU.
Quando a
\indexcode{swtch}
que estamos acompanhando retorna,
ela não volta para
\lstinline{sched},
mas sim para
\indexcode{scheduler},
com o ponteiro da pilha (stack pointer) apontando para a pilha do escalonador da CPU atual.

%%
\section{Code: Scheduling}
%%

A seção anterior examinou os detalhes de baixo nível da
\indexcode{swtch};
agora vamos considerar
\lstinline{swtch}
como um dado e analisar
a troca de uma kernel thread de um processo
para outro processo por meio do escalonador.
O escalonador existe na forma de uma thread especial por CPU, cada uma executando a função
\lstinline{scheduler}.
Essa função é responsável por escolher qual processo será executado em seguida.
Um processo que deseja liberar a CPU deve
adquirir o seu próprio lock de processo
\indexcode{p->lock},
liberar qualquer outro lock que esteja segurando,
atualizar seu próprio estado
(\lstinline{p->state}),
e então chamar
\indexcode{sched}.
Você pode ver essa sequência em
\lstinline{yield}
\lineref{kernel/proc.c:/^yield/},
\texttt{sleep}
e
\texttt{exit}.
\lstinline{sched}
verifica novamente alguns desses requisitos
\linerefs{kernel/proc.c:/if..holding/,/running/}
e então checa uma implicação:
como um lock está sendo mantido, as interrupções devem estar desabilitadas.
Finalmente,
\indexcode{sched}
chama
\indexcode{swtch}
para salvar o contexto atual em
\lstinline{p->context}
e trocar para o contexto do escalonador em
\indexcode{cpu->context}.
\lstinline{swtch}
retorna na pilha do escalonador
como se o
\indexcode{scheduler}'s
\lstinline{swtch}
tivesse retornado
\lineref{kernel/proc.c:/swtch.*contex.*contex/}.
O escalonador continua seu loop
\lstinline{for}, encontra um processo para rodar,
faz a troca para ele, e o ciclo se repete.

Acabamos de ver que o xv6 mantém
\indexcode{p->lock}
durante chamadas para
\lstinline{swtch}:
quem chama
\indexcode{swtch}
já deve estar com o lock adquirido, e o controle desse lock passa para o código que será executado após a troca.
Esse arranjo é incomum: é mais comum que a thread que adquire um lock também seja responsável por liberá-lo.
A troca de contexto no xv6 precisa quebrar essa convenção porque
\indexcode{p->lock}
protege invariantes sobre os campos
\lstinline{state}
e
\lstinline{context}
do processo que não são garantidos enquanto está se executando dentro de
\lstinline{swtch}.
Por exemplo, se
\lstinline{p->lock}
não estivesse mantido durante a
\indexcode{swtch},
uma CPU diferente poderia decidir executar o processo depois que
\indexcode{yield}
define seu estado como
\lstinline{RUNNABLE},
mas antes de
\lstinline{swtch}
fazer com que ele parasse de usar sua própria pilha do kernel.
O resultado seria duas CPUs usando a mesma pilha — o que causaria caos.
Uma vez que
\lstinline{yield}
começa a modificar o estado de um processo que está rodando para torná-lo
\lstinline{RUNNABLE},
\lstinline{p->lock}
deve permanecer adquirido até que os invariantes sejam restaurados:
o ponto mais cedo em que é seguro liberar o lock é após o
\lstinline{scheduler}
(em sua própria pilha) limpar
\lstinline{c->proc}.
Da mesma forma, uma vez que o
\lstinline{scheduler}
começa a converter um processo
\lstinline{RUNNABLE}
para
\lstinline{RUNNING},
o lock não pode ser liberado até que a kernel thread do processo esteja completamente em execução (após a
\lstinline{swtch},
por exemplo em
\lstinline{yield}).

O único lugar em que uma kernel thread libera sua CPU é em
\lstinline{sched},
e ela sempre faz a troca para o mesmo local em
\lstinline{scheduler}, que
(quase) sempre troca para alguma kernel thread que anteriormente chamou
\lstinline{sched}.
Assim, se alguém imprimisse os números das linhas onde o xv6 realiza trocas de threads, veria o seguinte padrão simples:
\lineref{kernel/proc.c:/swtch..c/},
\lineref{kernel/proc.c:/swtch..p/},
\lineref{kernel/proc.c:/swtch..c/},
\lineref{kernel/proc.c:/swtch..p/},
e assim por diante.
Procedimentos que transferem controle intencionalmente entre si por meio de troca de threads são às vezes chamados de
\indextext{coroutines};
neste exemplo,
\indexcode{sched}
e
\indexcode{scheduler}
são corrotinas um do outro.

Há um caso em que a chamada do escalonador para
\indexcode{swtch}
não termina em
\indexcode{sched}.
\lstinline{allocproc}
define o registrador
\lstinline{ra}
do contexto de um novo processo para
\indexcode{forkret}
\lineref{kernel/proc.c:/^forkret/},
de forma que sua primeira
\lstinline{swtch}
"retorna" para o início dessa função.
\lstinline{forkret}
existe para liberar o
\indexcode{p->lock};
caso contrário, como o novo processo precisa retornar ao espaço do usuário como se estivesse retornando de
\lstinline{fork},
ele poderia, ao invés disso, começar em
\lstinline{usertrapret}.
\lstinline{scheduler}
\lineref{kernel/proc.c:/^scheduler/}
executa um loop:
encontra um processo para rodar, executa até que ele chame \texttt{yield}, repete.
O escalonador percorre a tabela de processos procurando por um processo pronto para executar, aquele que tem
\lstinline{p->state}
\lstinline{==}
\lstinline{RUNNABLE}.
Assim que encontra um processo, ele define a variável de processo atual por CPU
\lstinline{c->proc},
marca o processo como
\lstinline{RUNNING},
e então chama
\indexcode{swtch}
para começar a executá-lo
\linerefs{kernel/proc.c:/Switch.to/,/swtch/}.

%% 
\section{Code: mycpu and myproc}
%% 
O xv6 frequentemente precisa de um ponteiro para a estrutura \lstinline{proc} do processo atual.
Em um uniprocessador, poderia haver uma variável global apontando para o \lstinline{proc} atual.
Isso não funciona em uma máquina multi-core, já que cada CPU executa um processo diferente.
A solução para esse problema é explorar o fato de que cada CPU tem seu próprio conjunto de registradores.

Enquanto uma CPU está executando no kernel, o xv6 garante que o registrador \lstinline{tp} da CPU sempre armazene o \textit{hartid} da CPU.
O RISC-V numera suas CPUs, dando a cada uma um \indextext{hartid} único.
\indexcode{mycpu}
\lineref{kernel/proc.c:/^mycpu/}
usa o \lstinline{tp} para indexar um array de estruturas \lstinline{cpu} e retornar aquela correspondente à CPU atual.
Uma
\indexcode{struct cpu}
\lineref{kernel/proc.h:/^struct.cpu/}
contém um ponteiro para a estrutura \lstinline{proc} do processo que está atualmente em execução naquela CPU (se houver),
registradores salvos para a thread do escalonador da CPU,
e o contador de spinlocks aninhados necessário para gerenciar o desligamento de interrupções.

Garantir que o \lstinline{tp} da CPU contenha o \textit{hartid} da CPU é um pouco mais complicado, pois o código de usuário pode modificar livremente o \lstinline{tp}.
A função \lstinline{start} define o registrador \lstinline{tp} no início da sequência de boot da CPU, enquanto ainda está no modo máquina
\lineref{kernel/start.c:/w_tp/}.
\lstinline{usertrapret} salva o \lstinline{tp} na página de trampolim, caso o código de usuário o modifique.
Por fim, \lstinline{uservec} restaura o \lstinline{tp} salvo ao entrar no kernel vindo do espaço do usuário
\lineref{kernel/trampoline.S:/make tp hold/}.
O compilador garante nunca modificar o \lstinline{tp} no código do kernel.
Seria mais conveniente se o xv6 pudesse consultar diretamente do hardware RISC-V o \textit{hartid} atual sempre que necessário,
mas o RISC-V permite isso apenas no modo máquina, não no modo supervisor.

Os valores retornados por
\lstinline{cpuid}
e
\lstinline{mycpu}
são frágeis: se o temporizador interromper e fizer com que a thread atual ceda a CPU e mais tarde retorne em outra CPU,
um valor retornado anteriormente não será mais correto.
Para evitar esse problema, o xv6 exige que os chamadores desabilitem as interrupções e só as reativem depois de terminar de usar a
\lstinline{struct cpu} retornada.

A função
\indexcode{myproc}
\lineref{kernel/proc.c:/^myproc/}
retorna o ponteiro para a
\lstinline{struct proc}
do processo que está rodando na CPU atual.
\lstinline{myproc}
desativa as interrupções, invoca
\lstinline{mycpu},
pega o ponteiro do processo atual
(\lstinline{c->proc})
a partir da
\lstinline{struct cpu},
e então reativa as interrupções.
O valor retornado por
\lstinline{myproc}
é seguro para usar mesmo que as interrupções estejam habilitadas:
se uma interrupção de timer mover o processo chamador para uma CPU diferente, seu ponteiro de
\lstinline{struct proc}
permanecerá o mesmo.

%% 
\section{Sleep and wakeup}
\label{sec:sleep}
%% 

O agendamento e os locks ajudam a esconder as ações de uma thread das outras,
mas também precisamos de abstrações que ajudem as threads a interagirem intencionalmente.
Por exemplo, o leitor de um pipe no xv6 pode precisar esperar que um processo escritor produza dados;
uma chamada \lstinline{wait} de um pai pode precisar esperar que o filho termine;
e um processo lendo o disco precisa esperar que o hardware do disco termine a leitura.
O kernel do xv6 usa um mecanismo chamado \textit{sleep} e \textit{wakeup} nessas (e muitas outras) situações.
\textit{Sleep} permite que uma kernel thread espere por um evento específico;
outra thread pode chamar \textit{wakeup} para indicar que threads esperando por um evento específico devem continuar.
\textit{Sleep} e \textit{wakeup} são frequentemente chamados de
\indextext{sequence coordination}
ou
\indextext{conditional synchronization}.

\textit{Sleep} e \textit{wakeup} oferecem uma interface de sincronização de nível relativamente baixo.
Para motivar a forma como eles funcionam no xv6, vamos usá-los para construir uma abstração de nível mais alto chamada
\indextext{semaphore}~\cite{dijkstra65},
que coordena produtores e consumidores
(o xv6 em si não usa semáforos).
Um semáforo mantém uma contagem e fornece duas operações.
A operação “V” (do produtor) incrementa a contagem.
A operação “P” (do consumidor) espera até que a contagem seja diferente de zero,
então a decrementa e retorna.
Se houvesse apenas uma thread produtora e uma consumidora,
e elas fossem executadas em CPUs diferentes,
e o compilador não otimizasse agressivamente,
essa implementação seria correta:
\begin{lstlisting}[numbers=left,firstnumber=100]
  struct semaphore {
    struct spinlock lock;
    int count;
  };

  void
  V(struct semaphore *s)
  {
     acquire(&s->lock);
     s->count += 1;
     release(&s->lock);
  }

  void
  P(struct semaphore *s)
  {
     while(s->count == 0)
       ;
     acquire(&s->lock);
     s->count -= 1;
     release(&s->lock);
  }
\end{lstlisting}
A implementação acima é custosa.
Se o produtor agir raramente, o consumidor passará a maior parte do tempo girando no laço
\lstinline{while}
na esperança de que a contagem se torne diferente de zero.
A CPU do consumidor provavelmente poderia realizar um trabalho mais produtivo do que
\indextext{espera ocupada (busy waiting)}
verificando repetidamente
\indextext{(polling)}
\lstinline{s->count}.
Evitar espera ocupada requer uma forma de o consumidor ceder a CPU
e retomar apenas após
\lstinline{V}
incrementar a contagem.

Aqui está um passo nessa direção, embora, como veremos, isso ainda não é suficiente.
Vamos imaginar um par de chamadas,
\indexcode{sleep}
e
\indexcode{wakeup},
que funcionam da seguinte maneira:
\lstinline{sleep(chan)}
espera por um evento designado pelo valor de
\indexcode{chan},
chamado de
\indextext{canal de espera (wait channel)}.
\lstinline{sleep}
coloca o processo que está chamando para dormir, liberando a CPU
para que outros processos possam utilizá-la.
\lstinline{wakeup(chan)}
acorda todos os processos que estão em chamadas
para \lstinline{sleep} com o mesmo
\lstinline{chan}
(se houver), fazendo com que suas chamadas para
\lstinline{sleep}
retornem.
Se nenhum processo estiver esperando em
\lstinline{chan},
\lstinline{wakeup}
não faz nada.
Podemos alterar a implementação do semáforo para usar
\lstinline{sleep}
e
\lstinline{wakeup} (alterações destacadas em amarelo):

c
Copiar
Editar
void
V(struct semaphore *s)
{
   acquire(&s->lock);
   s->count += 1;
   wakeup(s);
   release(&s->lock);
}

void
P(struct semaphore *s)
{
  while(s->count == 0)
    sleep(s);
  acquire(&s->lock);
  s->count -= 1;
  release(&s->lock);
}
\end{lstlisting}
% \begin{figure}[t]
% \center
% \includegraphics[scale=0.5]{fig/deadlock.pdf}
% \caption{Example lost wakeup problem}
% \label{fig:deadlock}
% \end{figure}
\lstinline{P}
agora libera a CPU em vez de ficar girando em espera, o que é uma melhoria.
No entanto, não é simples projetar
\lstinline{sleep}
e
\lstinline{wakeup}
com essa interface sem sofrer com o que é conhecido como o problema do
\indextext{lost wake-up} (sinal perdido).

Suponha que
\lstinline{P}
descubra que
\lstinline{s->count == 0}
na linha marcada como test.
Enquanto
\lstinline{P}
está entre as linhas test e sleep,
\lstinline{V}
executa em outra CPU:
ela muda
\lstinline{s->count}
para um valor diferente de zero e chama
\lstinline{wakeup},
que não encontra nenhum processo dormindo e, portanto, não faz nada.
Agora
\lstinline{P}
continua sua execução na linha sleep:
chama
\lstinline{sleep}
e entra em modo de espera.
Isso causa um problema:
\lstinline{P}
está dormindo esperando por uma chamada de
\lstinline{V}
que já aconteceu.
A menos que tenhamos sorte e o produtor chame
\lstinline{V}
novamente, o consumidor esperará para sempre,
mesmo que a contagem seja diferente de zero.

A raiz do problema é que o invariante de que
\lstinline{P}
só dorme quando
\lstinline{s->count == 0}
é violado porque
\lstinline{V}
executa no momento exato errado.
Uma maneira incorreta de proteger esse invariante seria mover a aquisição do lock (destacada em amarelo abaixo) em
\lstinline{P}
de modo que a verificação da contagem e a chamada para
\lstinline{sleep}
sejam atômicas:
\begin{lstlisting}[numbers=left,firstnumber=300]
  void
  V(struct semaphore *s)
  {
    acquire(&s->lock);
    s->count += 1;
    wakeup(s);
    release(&s->lock);
  }
  
  void
  P(struct semaphore *s)
  {
    (*@\hl{acquire(\&s->lock);}@*)
    while(s->count == 0)    (*@\label{line:test1}@*)
      sleep(s);             (*@\label{line:sleep1}@*)
    s->count -= 1;
    release(&s->lock);
  }
\end{lstlisting}
Pode-se esperar que essa versão de
\lstinline{P}
evite o problema de wakeup perdido porque o lock impede
\lstinline{V}
de ser executado entre as linhas~\ref{line:test1} e~\ref{line:sleep1}.
De fato, isso acontece, mas também gera um deadlock:
\lstinline{P}
mantém o lock enquanto dorme, então
\lstinline{V}
ficará bloqueado para sempre esperando pelo lock.

Vamos corrigir esse esquema anterior alterando a interface da
\lstinline{sleep}:
o chamador deve passar o \indextext{condition lock} para
\lstinline{sleep},
para que ela possa liberar o lock depois que o processo chamador for marcado como adormecido e estiver esperando no canal de sleep.
O lock forçará um
\lstinline{V}
concorrente a esperar até que
\lstinline{P}
tenha terminado de se colocar para dormir, garantindo que o
\lstinline{wakeup}
encontre o consumidor adormecido e o acorde.
Assim que o consumidor estiver acordado novamente,
\indexcode{sleep}
readquire o lock antes de retornar.
Nosso novo e correto esquema de sleep/wakeup pode ser usado da seguinte forma (alteração destacada em amarelo):
\begin{lstlisting}[numbers=left,firstnumber=400]
  void
  V(struct semaphore *s)
  {
    acquire(&s->lock);
    s->count += 1;
    wakeup(s);
    release(&s->lock);
  }

  void
  P(struct semaphore *s)
  {
    acquire(&s->lock);
    while(s->count == 0)
       (*@\hl{sleep(s, \&s->lock);}@*)
    s->count -= 1;
    release(&s->lock);
  }
\end{lstlisting}
O fato de que
\lstinline{P}
mantém
\lstinline{s->lock}
impede que
\lstinline{V}
tente acordá-lo entre a verificação de
\lstinline{s->count}
por parte de
\lstinline{P}
e sua chamada para
\lstinline{sleep}.
No entanto,
\lstinline{sleep}
deve liberar
\lstinline{s->lock}
e colocar o processo consumidor para dormir
de forma atômica do ponto de vista do
\lstinline{wakeup},
a fim de evitar perdas de notificações (lost wakeups).
%% 
\section{Code: Sleep and wakeup}
%% 
As funções
\indexcode{sleep}
\lineref{kernel/proc.c:/^sleep/}
e
\indexcode{wakeup}
\lineref{kernel/proc.c:/^wakeup/}
do Xv6 fornecem a interface usada no último exemplo acima.
A ideia básica é que
\lstinline{sleep}
marque o processo atual como
\indexcode{SLEEPING}
e então chame
\indexcode{sched}
para liberar a CPU;
\lstinline{wakeup}
procura um processo dormindo no canal de espera fornecido
e o marca como
\indexcode{RUNNABLE}.
Quem chama
\lstinline{sleep}
e
\lstinline{wakeup}
pode usar qualquer número mutuamente conveniente como canal.
O Xv6 frequentemente usa o endereço
de alguma estrutura de dados do kernel envolvida na espera.

\lstinline{sleep}
adquire primeiro
\indexcode{p->lock}
\lineref{kernel/proc.c:/DOC: sleeplock1/}
e só então libera
\lstinline{lk}.
Como veremos, o fato de que
\lstinline{sleep}
mantém um dos dois locks (ou ambos) o tempo todo
é o que impede um
\lstinline{wakeup}
concorrente (que deve adquirir e manter ambos) de agir.
Agora que
\lstinline{sleep}
mantém apenas
\lstinline{p->lock},
ele pode colocar o processo para dormir registrando
o canal de sleep,
alterando o estado do processo para \texttt{SLEEPING},
e chamando
\lstinline{sched}
\linerefs{kernel/proc.c:/chan.=.chan/,/sched/}.
Em breve ficará claro por que é crucial que
\lstinline{p->lock}
não seja liberado (pelo
\lstinline{scheduler}) até que o processo esteja marcado como \texttt{SLEEPING}.
Em algum momento, um processo irá adquirir o lock de condição,
configurar a condição pela qual o processo adormecido está esperando
e chamar \lstinline{wakeup(chan)}.
É importante que \lstinline{wakeup} seja chamado
enquanto o lock de condição estiver mantido\footnote{% % Tecnicamente, é suficiente que \lstinline{wakeup}
ocorra após o \lstinline{acquire}
(ou seja, pode-se chamar \lstinline{wakeup}
depois do \lstinline{release}).% % }.
\lstinline{wakeup} percorre a tabela de processos
\lineref{kernel/proc.c:/^wakeup(/}.
Ela adquire o \lstinline{p->lock}
de cada processo que inspeciona.
Quando \lstinline{wakeup} encontra um processo no estado
\indexcode{SLEEPING}
com um \indexcode{chan} correspondente,
ela altera o estado desse processo para
\indexcode{RUNNABLE}.
Na próxima vez que o \lstinline{scheduler} for executado,
ele verá que o processo está pronto para rodar.

Por que as regras de locking para
\lstinline{sleep}
e
\lstinline{wakeup}
garantem que um processo que vai dormir
não perderá um wakeup concorrente?
O processo que está indo dormir mantém
ou o lock de condição ou seu próprio
\lstinline{p->lock}
ou ambos desde antes de verificar a condição
até depois de ser marcado como \texttt{SLEEPING}.
O processo que chama \texttt{wakeup} mantém ambos os locks
no loop da função \texttt{wakeup}.
Assim, o processo que acorda ou
torna a condição verdadeira antes que a thread consumidora a verifique;
ou a \lstinline{wakeup} da thread acordadora examina a thread adormecida
estritamente depois que ela foi marcada como \texttt{SLEEPING}.
Nesse caso, \lstinline{wakeup} verá o processo adormecido
e o acordará (a menos que algo mais o acorde antes).

Às vezes, múltiplos processos estão dormindo
no mesmo canal; por exemplo, mais de um processo
lendo de um pipe.
Uma única chamada para
\lstinline{wakeup}
acordará todos eles.
Um deles será executado primeiro e adquirirá o lock com o qual
\lstinline{sleep} foi chamado, e (no caso de pipes) lerá os dados
que estiverem disponíveis.
Os outros processos descobrirão que, apesar de terem sido acordados,
não há dados para serem lidos.
Do ponto de vista deles, o wakeup foi “espúrio”, e
eles precisarão dormir novamente.
Por esse motivo, \lstinline{sleep} é sempre chamado dentro de um loop que
verifica a condição.
Nenhum dano é causado se dois usos de sleep/wakeup acidentalmente
escolherem o mesmo canal: eles verão acordamentos espúrios,
mas o loop descrito anteriormente tolera esse problema.
Grande parte da elegância de sleep/wakeup está no fato de que
ele é leve (não é necessário criar estruturas de dados especiais
para atuar como canais de espera) e fornece uma camada de indireção
(os chamadores não precisam saber com qual processo específico
estão interagindo).

%% \section{Código: Pipes} %%

Um exemplo mais complexo que usa \lstinline{sleep}
e \lstinline{wakeup} para sincronizar produtores e
consumidores é a implementação de pipes no xv6.
Vimos a interface de pipes no Capítulo~\ref{CH:UNIX}:
bytes escritos em uma extremidade de um pipe são copiados
para um buffer dentro do kernel e podem então ser lidos
pela outra extremidade do pipe.
Capítulos futuros examinarão o suporte a descritores de arquivo
em torno dos pipes, mas vamos olhar agora para as
implementações de
\indexcode{pipewrite}
e
\indexcode{piperead}.

Cada pipe
é representado por uma
\indexcode{struct pipe},
que contém
um \lstinline{lock}
e um buffer de
\lstinline{data}.
Os campos
\lstinline{nread}
e
\lstinline{nwrite}
contam o número total de bytes lidos e escritos no buffer.
O buffer é circular:
o próximo byte escrito após
\lstinline{buf[PIPESIZE-1]}
é
\lstinline{buf[0]}.
Essas contagens não se repetem (não fazem wrap).
Essa convenção permite que a implementação
diferencie um buffer cheio
(\lstinline{nwrite}
\lstinline{==}
\lstinline{nread+PIPESIZE})
de um buffer vazio
(\lstinline{nwrite}
\lstinline{==}
\lstinline{nread}),
mas isso significa que o índice no buffer
deve ser calculado usando
\lstinline{buf[nread % PIPESIZE]}
ao invés de apenas
\lstinline{buf[nread]}
(e de forma semelhante para
\lstinline{nwrite}).

Vamos supor que chamadas para
\lstinline{piperead}
e
\lstinline{pipewrite}
ocorram simultaneamente em dois CPUs diferentes.
\lstinline{pipewrite}
\lineref{kernel/pipe.c:/^pipewrite/}
começa adquirindo o lock do pipe, o qual
protege as contagens, os dados e seus
invariantes associados.
\lstinline{piperead}
\lineref{kernel/pipe.c:/^piperead/}
então tenta adquirir o mesmo lock, mas não consegue.
Ele fica em espera no
\lstinline{acquire}
\lineref{kernel/spinlock.c:/^acquire/}
aguardando pelo lock.
Enquanto
\lstinline{piperead}
espera,
\lstinline{pipewrite}
faz um loop sobre os bytes sendo escritos
(\lstinline{addr[0..n-1]}),
adicionando cada um ao pipe
\lineref{kernel/pipe.c:/nwrite++/}.
Durante esse loop, pode acontecer de
o buffer encher
\lineref{kernel/pipe.c:/DOC: pipewrite-full/}.
Nesse caso,
\lstinline{pipewrite}
chama
\lstinline{wakeup}
para alertar quaisquer leitores adormecidos de que
há dados esperando no buffer
e então dorme sobre
\lstinline{&pi->nwrite}
para aguardar que um leitor retire alguns bytes
do buffer.
\lstinline{sleep}
libera
o lock do pipe
como parte do processo de colocar
o processo de
\lstinline{pipewrite}
para dormir.
\lstinline{piperead}
agora adquire o lock do pipe e entra em sua seção crítica:
ele verifica que
\lstinline{pi->nread}
\lstinline{!=}
\lstinline{pi->nwrite}
\lineref{kernel/pipe.c:/DOC: pipe-empty/}
(\lstinline{pipewrite}
foi colocado para dormir porque
\lstinline{pi->nwrite}
\lstinline{==}
\lstinline{pi->nread}
\lstinline{+}
\lstinline{PIPESIZE}
\lineref{kernel/pipe.c:/pipewrite-full/}),
então ele prossegue para o loop
\lstinline{for}, copia dados do pipe
\lineref{kernel/pipe.c:/DOC: piperead-copy/},
e incrementa
\lstinline{nread}
pela quantidade de bytes copiados.
Essa quantidade de bytes agora está disponível para escrita, então
\lstinline{piperead}
chama
\lstinline{wakeup}
\lineref{kernel/pipe.c:/DOC: piperead-wakeup/}
para acordar quaisquer escritores adormecidos
antes de retornar.
\lstinline{wakeup}
encontra um processo dormindo em
\lstinline{&pi->nwrite},
o processo que estava executando
\lstinline{pipewrite}
mas foi interrompido quando o buffer encheu.
Ele marca esse processo como
\indexcode{RUNNABLE}.

O código do pipe usa canais de sono separados para leitores e escritores
(\lstinline{pi->nread}
e
\lstinline{pi->nwrite});
isso pode tornar o sistema mais eficiente no caso improvável
de haver muitos
leitores e escritores esperando pelo mesmo pipe.
O código do pipe dorme dentro de um loop que verifica a
condição de sono; se houver múltiplos leitores
ou escritores, todos, exceto o primeiro processo a acordar,
verão que a condição ainda é falsa e voltarão a dormir.

%% \section{Código: Wait, exit, and kill} %%

\lstinline{sleep}
e
\lstinline{wakeup}
podem ser usados para vários tipos de espera.
Um exemplo interessante, introduzido no Capítulo~\ref{CH:UNIX},
é a interação entre o \indexcode{exit} de um filho
e o \indexcode{wait} de seu pai.
No momento da morte do filho, o pai pode já
estar dormindo em {\tt wait}, ou pode estar fazendo outra coisa;
neste último caso, uma chamada subsequente a {\tt wait} deve
observar a morte do filho, talvez muito tempo depois de ele chamar {\tt exit}.
A maneira como o xv6 registra a morte do filho até que {\tt wait}
a observe é fazer com que {\tt exit} coloque o processo chamador no estado \indexcode{ZOMBIE},
onde ele permanece até que o {\tt wait} do pai o detecte, mude
o estado do filho para {\tt UNUSED}, copie o status de saída do filho,
e retorne o ID do processo do filho para o pai.
Se o pai terminar antes do filho, o
pai entrega o filho ao processo
\lstinline{init},
que chama {\tt wait} continuamente;
assim,
todo filho tem um pai para limpar seus recursos.
Um desafio é
evitar condições de corrida e deadlocks entre
pai e filho simultâneos.
\lstinline{wait}
e
\lstinline{exit},
assim como chamadas simultâneas de
\lstinline{exit} e \lstinline{exit}.

\lstinline{wait} começa adquirindo
\lstinline{wait_lock}
\lineref{kernel/proc.c:/^wait/},
que atua como o lock de condição que ajuda a garantir que \lstinline{wait} não perca um \lstinline{wakeup}
de um filho que está saindo.
Em seguida, \lstinline{wait} percorre a tabela de processos.
Se encontrar um filho no estado \texttt{ZOMBIE},
ele libera os recursos desse filho e
sua estrutura \lstinline{proc}, copia
o status de saída do filho para o endereço fornecido a \lstinline{wait}
(se não for 0),
e retorna o ID do processo do filho.
Se
\lstinline{wait}
encontra filhos mas nenhum saiu ainda,
ele chama
\lstinline{sleep}
para esperar que algum deles saia
\lineref{kernel/proc.c:/DOC: wait-sleep/},
e então escaneia novamente.
\lstinline{wait} frequentemente mantém dois locks:
\lstinline{wait_lock} e o \lstinline{pp->lock} de algum processo;
a ordem para evitar deadlocks é primeiro \lstinline{wait_lock}
e depois \lstinline{pp->lock}.

\lstinline{exit} \lineref{kernel/proc.c:/^exit/} registra o status
de saída, libera alguns recursos, chama \lstinline{reparent} para passar seus
filhos ao processo \lstinline{init}, acorda o pai caso
ele esteja em \lstinline{wait}, marca o chamador como um zumbi
e cede permanentemente a CPU. \lstinline{exit} mantém tanto
\lstinline{wait_lock} quanto \lstinline{p->lock} durante essa
sequência.
Ele mantém \lstinline{wait_lock} porque
é o lock de condição para o \lstinline{wakeup(p->parent)}, evitando que um pai em
\lstinline{wait} perca o wakeup. \lstinline{exit} também deve manter
\lstinline{p->lock} durante essa sequência para evitar que o pai em
\lstinline{wait} veja que o filho está no estado
\lstinline{ZOMBIE} antes que o filho tenha finalmente chamado
\lstinline{swtch}. \lstinline{exit} adquire esses locks na
mesma ordem que \lstinline{wait} para evitar deadlocks.

Pode parecer incorreto que \lstinline{exit} acorde o pai
antes de definir seu estado como \lstinline{ZOMBIE},
mas isso é seguro:
embora
\lstinline{wakeup}
possa fazer com que o pai seja executado,
o loop em
\lstinline{wait}
não pode examinar o filho até que o
\lstinline{p->lock} do filho seja liberado pelo {\tt scheduler},
então
\lstinline{wait}
não pode observar
o processo que está saindo até bem depois de
\lstinline{exit}
ter definido seu estado como
\lstinline{ZOMBIE}
\lineref{kernel/proc.c:/state.=.ZOMBIE/}.

Enquanto
\lstinline{exit}
permite que um processo termine a si mesmo,
\lstinline{kill}
\lineref{kernel/proc.c:/^kill/}
permite que um processo solicite que outro termine.
Seria muito complexo para
\lstinline{kill}
destruir diretamente o processo alvo, pois o alvo
poderia estar executando em outra CPU, talvez
no meio de uma sequência sensível de atualizações em estruturas de dados do kernel.
Assim,
\lstinline{kill}
faz muito pouco: apenas define
\indexcode{p->killed}
do processo alvo e, se ele estiver dormindo, o acorda.
Eventualmente o alvo entrará ou sairá do kernel,
momento no qual o código em
\lstinline{usertrap}
chamará
\lstinline{exit}
se
\lstinline{p->killed}
estiver definido
(ele verifica isso chamando
\lstinline{killed}
\lineref{kernel/proc.c:/^killed/}).
Se o alvo estiver rodando em espaço de usuário, ele em breve entrará
no kernel fazendo uma chamada de sistema ou porque o temporizador (ou
algum outro dispositivo) o interrompeu.

Se o processo vítima estiver em
\lstinline{sleep},
a chamada de
\lstinline{kill}
para
\lstinline{wakeup}
fará com que o processo vítima retorne de
\lstinline{sleep}.
Isso é potencialmente perigoso porque
a condição pela qual ele estava esperando pode não ser mais verdadeira.
No entanto, chamadas para
\lstinline{sleep}
em xv6 estão sempre envolvidas em um
\lstinline{while}
que reavalia a condição após o retorno de
\lstinline{sleep}.
Algumas chamadas para
\lstinline{sleep}
também verificam
\lstinline{p->killed}
no loop e abandonam a atividade atual se esse campo estiver ativado.
Isso só é feito quando esse abandono é apropriado.
Por exemplo, o código de leitura e escrita de pipe
\lineref{kernel/pipe.c:/killed.pr/}
retorna se a flag killed estiver definida; eventualmente o
código retornará ao trap, que novamente
verificará \lstinline{p->killed} e chamará \lstinline{exit}.

Alguns loops de
\lstinline{sleep}
no xv6 não verificam
\lstinline{p->killed}
porque o código está no meio de uma chamada de sistema
com múltiplas etapas que devem ser atômicas.
O driver virtio
\lineref{kernel/virtio_disk.c:/sleep.b/}
é um exemplo: ele não verifica
\lstinline{p->killed}
porque uma operação de disco pode ser uma entre várias
escritas necessárias para manter o sistema de arquivos
em um estado consistente.
Um processo que for morto enquanto espera por uma operação de disco
não irá sair imediatamente; ele completará a chamada de sistema atual
e então \lstinline{usertrap} verá a flag killed.

\section{Process Locking}
O lock associado a cada processo (\lstinline{p->lock}) é o
lock mais complexo em xv6.
Uma forma simples de entender \lstinline{p->lock} é
que ele deve ser mantido enquanto se lê ou escreve qualquer
um dos seguintes campos da estrutura \lstinline{struct proc}:
\lstinline{p->state},
\lstinline{p->chan},
\lstinline{p->killed},
\lstinline{p->xstate}
e
\lstinline{p->pid}.
Esses campos podem ser acessados por outros processos ou por
threads do escalonador em outras CPUs, então é natural
que precisem ser protegidos por um lock.

No entanto, a maioria dos usos de \lstinline{p->lock} está relacionada
à proteção de aspectos de mais alto nível das estruturas de dados
e algoritmos de processo do xv6. Aqui está o conjunto completo
de situações em que \lstinline{p->lock}

O campo \lstinline{p->parent} é protegido pelo lock global
\lstinline{wait_lock}, e não por \lstinline{p->lock}.
Apenas o processo pai modifica \lstinline{p->parent},
embora o campo seja lido tanto pelo próprio processo quanto por
outros que procuram por seus filhos.
O propósito de
\lstinline{wait_lock}
é atuar como lock de condição quando
\lstinline{wait}
entra em espera aguardando que algum filho termine.
Um filho que está saindo mantém ou \lstinline{wait_lock}
ou \lstinline{p->lock}.
até depois de definir seu estado como \lstinline{ZOMBIE}, acordar
seu processo pai e ceder a CPU. \lstinline{wait_lock} também
serializa chamadas concorrentes de \lstinline{exit} feitas por pai e filho,
de forma que o processo \lstinline{init} (que herda o filho) seja
garantidamente acordado de sua chamada de \lstinline{wait}.
\lstinline{wait_lock} é um lock global, e não um lock por processo em cada pai,
porque, até que um processo o adquira, ele não sabe quem é seu pai.

Mundo Real
O escalonador do xv6 implementa uma política de escalonamento simples,
que executa cada processo por sua vez. Essa política é chamada de
round robin.
Sistemas operacionais reais implementam políticas mais sofisticadas que, por exemplo,
permitem que processos tenham prioridades. A ideia é que um processo pronto de alta prioridade
tenha preferência sobre um processo pronto de baixa prioridade.
Essas políticas podem se tornar rapidamente complexas, porque há objetivos conflitantes:
por exemplo, o sistema operacional pode querer garantir justiça
e alta vazão ao mesmo tempo.
Além disso, políticas complexas podem levar a interações indesejadas, como
inversão de prioridade (priority inversion)
e
comboios (convoys).
A inversão de prioridade pode ocorrer quando processos de alta e baixa prioridade
compartilham um mesmo lock, e o processo de baixa prioridade impede o de alta
de continuar sua execução por ter adquirido o lock.
Um longo comboio de processos pode se formar quando vários processos de alta prioridade
ficam esperando um de baixa que está com um lock compartilhado; uma vez que esse comboio se forma,
ele pode persistir por um longo tempo.
Para evitar esses problemas, escalonadores sofisticados precisam de mecanismos adicionais.

\lstinline{sleep}
e
\lstinline{wakeup}
são métodos simples e eficazes de sincronização,
mas existem muitos outros.
O primeiro desafio de todos eles é
evitar o problema de "sinal perdido" (lost wakeups) que vimos no
início do capítulo.
O kernel original do Unix
\lstinline{sleep}
simplesmente desabilitava interrupções,
o que era suficiente porque Unix rodava em um sistema de CPU única.
Como o xv6 roda em multiprocessadores,
ele adiciona um lock explícito ao
\lstinline{sleep}.
O \lstinline{msleep} do FreeBSD
adota a mesma abordagem.
O \lstinline{sleep} do Plan 9
usa uma função de callback que é executada com o lock do escalonador
mantido, pouco antes de dormir;
essa função faz uma verificação de última hora
da condição de sono, para evitar sinalizações perdidas.
O \lstinline{sleep} do kernel do Linux
usa uma fila de processos explícita, chamada wait queue, em vez de
um canal de espera (wait channel); a fila possui seu próprio lock interno.

Vasculhar todos os processos em
\lstinline{wakeup}
é ineficiente. Uma solução melhor é
substituir o
\lstinline{chan}
tanto em
\lstinline{sleep}
quanto em
\lstinline{wakeup}
por uma estrutura de dados que mantém
uma lista de processos dormindo sobre essa estrutura,
como a wait queue do Linux.
O \lstinline{sleep} e \lstinline{wakeup} do Plan 9
chamam essa estrutura de ponto de encontro (rendezvous point).
Muitas bibliotecas de threads se referem a essa mesma
estrutura como uma variável de condição (condition variable);
nesse contexto, as operações
\lstinline{sleep}
e
\lstinline{wakeup}
são chamadas de
\lstinline{wait}
e
\lstinline{signal}.
Todos esses mecanismos compartilham uma ideia central:
a condição de sono é protegida por algum tipo de lock
que é liberado atomica-mente durante o \lstinline{sleep}.

A implementação de
\lstinline{wakeup}
acorda todos os processos que estão esperando em um determinado canal,
e pode ser o caso de muitos processos estarem esperando nesse canal.
O sistema operacional então escalona todos esses processos, e eles disputam
para verificar a condição de sono. Processos que agem dessa forma às vezes
são chamados de
thundering herd (manada trovejante),
e é melhor evitar isso.
A maioria das variáveis de condição possui dois tipos de operação para
\lstinline{wakeup}:
\lstinline{signal},
que acorda um processo, e
\lstinline{broadcast},
que acorda todos os processos esperando.

Semáforos são frequentemente usados para sincronização.
A contagem normalmente corresponde a algo como
o número de bytes disponíveis no buffer de um pipe
ou o número de filhos zumbis que um processo possui.
Usar uma contagem explícita como parte da abstração
evita o problema de "sinal perdido":
há uma contagem explícita de quantos
\lstinline{wakeups}
ocorreram.
A contagem também evita problemas como sinalizações espúrias
e a formação de manadas (\lstinline{thundering herd}).

Encerrar processos e liberar seus recursos introduz muita complexidade no xv6.
Na maioria dos sistemas operacionais, isso é ainda mais complexo, pois, por exemplo,
o processo vítima pode estar profundamente dentro do kernel dormindo, e desfazer
sua pilha de chamadas exige cuidado, já que cada função na pilha
pode precisar realizar uma limpeza.
Algumas linguagens ajudam com mecanismos de exceção, mas C não.
Além disso, há outros eventos que podem acordar um processo que está dormindo,
mesmo que o evento esperado ainda não tenha ocorrido. Por exemplo, quando um processo
Unix está dormindo, outro processo pode enviar um
signal
a ele. Nesse caso, o processo retorna da chamada de sistema interrompida com valor -1
e com o código de erro definido como EINTR.
A aplicação pode verificar esses valores e decidir o que fazer.
O xv6 não suporta sinais, então essa complexidade não ocorre.

O suporte do xv6 a
\lstinline{kill}
não é totalmente satisfatório: existem loops de sono
que provavelmente deveriam verificar
\lstinline{p->killed}.
Um problema relacionado é que, mesmo para loops de
\lstinline{sleep}
que verificam
\lstinline{p->killed},
existe uma condição de corrida entre
\lstinline{sleep}
e
\lstinline{kill};
o segundo pode definir
\lstinline{p->killed}
e tentar acordar o processo vítima logo após o loop da vítima
verificar \lstinline{p->killed}
mas antes de chamar
\lstinline{sleep}.
Se isso ocorrer, a vítima não notará que foi marcada como
\lstinline{killed}
até que a condição pela qual esperava aconteça. Isso pode demorar bastante
ou talvez nunca ocorra
(por exemplo, se a vítima estiver esperando entrada do console, mas o usuário
não digitar nada).

Um sistema operacional real encontraria estruturas \lstinline{proc}
livres usando uma lista livre explícita
em tempo constante, em vez da busca linear usada em
\lstinline{allocproc};
o xv6 usa a busca linear por simplicidade.
\end{enumerate}
